\documentclass[UTF8]{ctexart}
\usepackage{amsmath, amssymb, geometry}
\usepackage{xcolor}
\usepackage{enumitem}
\usepackage{tcolorbox}
\usepackage{cases}
\usepackage{fancyhdr}
\geometry{a4paper, margin=1in}

% 定义红色环境
\newenvironment{redenv}{\color{red}}{}

% 定义详细解析环境
\newenvironment{detailedanalysis}{%
    \color{red}
    \begin{enumerate}[label=\textbf{第\arabic*步：}, leftmargin=*, align=left, itemsep=1.5ex]
}{%
    \end{enumerate}
}

% 定义概念框
\newenvironment{conceptbox}{%
    \begin{tcolorbox}[colback=red!5, colframe=red!50!black, title=概念解析, fonttitle=\bfseries]
}{%
    \end{tcolorbox}
}

% 定义公式推导环境
\newenvironment{formuladerivation}{%
    \begin{enumerate}[label={\textbf{公式\arabic*：}}, leftmargin=*, align=left]
}{%
    \end{enumerate}
}

% 定义题目和解析的分隔线
\newcommand{\problemrule}{\noindent\rule{\textwidth}{1.5pt}\vspace{-0.8em}}
\newcommand{\analysisrule}{\noindent\rule{\textwidth}{0.8pt}\vspace{-0.8em}}

% 宏定义：方便排版选择题选项
\newcommand{\choicesFour}[4]{
    \begin{enumerate}[label=\Alph*., itemsep=0pt, parsep=0pt, topsep=5pt, labelsep=0.5em]
        \begin{minipage}{0.22\linewidth} \item #1 \end{minipage}
        \begin{minipage}{0.22\linewidth} \item #2 \end{minipage}
        \begin{minipage}{0.22\linewidth} \item #3 \end{minipage}
        \begin{minipage}{0.22\linewidth} \item #4 \end{minipage}
    \end{enumerate}
}
\newcommand{\choicesTwo}[4]{
    \begin{enumerate}[label=\Alph*., itemsep=0pt, parsep=0pt, topsep=5pt, labelsep=0.5em]
        \begin{minipage}{0.45\linewidth} \item #1 \end{minipage}
        \begin{minipage}{0.45\linewidth} \item #2 \end{minipage}
        \begin{minipage}{0.45\linewidth} \item #3 \end{minipage}
        \begin{minipage}{0.45\linewidth} \item #4 \end{minipage}
    \end{enumerate}
}
\newcommand{\choices}[4]{
    \begin{enumerate}[label=\Alph*., itemsep=0pt, parsep=0pt, topsep=5pt, labelsep=0.5em]
        \item #1
        \item #2
        \item #3
        \item #4
    \end{enumerate}
}

\newcommand{\blank}[1]{\underline{\hspace{#1}}}

\begin{document}

\section*{第十章 向量代数与空间解析几何（提高）}

% ========== 第1题 ==========
\problemrule
\section*{【题目1】直线夹角计算}
直线 $ L_{1} $ : $ x - 1 = \frac{y - 5}{-2} = z + 8 $ 与直线 $ L_{2} $ : $ \begin{cases} x - y = 6 \\ 2y + z = 3 \end{cases} $ 的夹角为 \blank{1cm}。
\choicesFour{$ \frac{\pi}{6} $}{$ \frac{\pi}{4} $}{$ \frac{\pi}{3} $}{$ \frac{\pi}{2} $}

\begin{redenv}
\analysisrule
\subsection*{逐步解析}

\begin{detailedanalysis}
\item \textbf{详细求解过程}
\begin{enumerate}[label=(\arabic*)]
\item \textbf{确定 $L_1$ 方向向量 $\mathbf{s}_1$}：
$L_1$ 的标准式给出 $\mathbf{s}_1 = \{1, -2, 1\}$。
\item \textbf{确定 $L_2$ 方向向量 $\mathbf{s}_2$}：
$L_2$ 由两平面交线确定，方向向量为法向量的叉积。
$\mathbf{n}_1 = \{1, -1, 0\}$，$\mathbf{n}_2 = \{0, 2, 1\}$。
$\mathbf{s}_2 = \mathbf{n}_1 \times \mathbf{n}_2 = \begin{vmatrix} \mathbf{i} & \mathbf{j} & \mathbf{k} \\ 1 & -1 & 0 \\ 0 & 2 & 1 \end{vmatrix} = (-1-0)\mathbf{i} - (1-0)\mathbf{j} + (2-0)\mathbf{k} = \{-1, -1, 2\}$。
\item \textbf{计算夹角 $\theta$}：
$\cos \theta = \frac{|\mathbf{s}_1 \cdot \mathbf{s}_2|}{|\mathbf{s}_1| |\mathbf{s}_2|} = \frac{|1(-1) + (-2)(-1) + 1(2)|}{\sqrt{1+4+1}\sqrt{1+1+4}} = \frac{|-1+2+2|}{\sqrt{6}\sqrt{6}} = \frac{3}{6} = \frac{1}{2}$。
$\theta = \frac{\pi}{3}$。
\end{enumerate}

\item \textbf{最终答案}
\textbf{答案}：C
\end{detailedanalysis}
\end{redenv}

% ========== 第2题 ==========
\problemrule
\section*{【题目2】切平面方程}
曲面 $ z = 4 - x^{2} - y^{2} $ 上一点 $P$ 的切平面平行于平面 $ 2x + 2y + z - 1 = 0 $ ，则 $P$ 点坐标为 \blank{1cm}。
\choicesFour{$ (1,-1,2) $}{$ (-1,1,2) $}{$ (1,1,2) $}{$ (-1,-1,2) $}

\begin{redenv}
\analysisrule
\subsection*{逐步解析}

\begin{detailedanalysis}
\item \textbf{详细求解过程}
\begin{enumerate}[label=(\arabic*)]
\item \textbf{确定法向量}：
曲面方程可写为 $F(x,y,z) = x^2 + y^2 + z - 4 = 0$。
法向量 $\mathbf{n} = \{F_x, F_y, F_z\} = \{2x, 2y, 1\}$。
\item \textbf{平行条件}：
切平面平行于 $2x + 2y + z - 1 = 0$，其法向量 $\mathbf{n}_0 = \{2, 2, 1\}$。
平行意味着 $\mathbf{n}$ 与 $\mathbf{n}_0$ 共线。
因为 $z$ 分量都是1，所以直接相等：
$2x = 2 \implies x = 1$。
$2y = 2 \implies y = 1$。
\item \textbf{求 $z$ 坐标}：
$z = 4 - 1^2 - 1^2 = 2$。
点 $P(1, 1, 2)$。
\end{enumerate}

\item \textbf{最终答案}
\textbf{答案}：C
\end{detailedanalysis}
\end{redenv}

% ========== 第3题 ==========
\problemrule
\section*{【题目3】线面位置关系}
直线 $l$: $ \frac{x+3}{-2}=\frac{y+4}{-7}=\frac{z}{3} $ 与平面 $ \pi $: $4x-2y-2z-3=0$ 的位置关系是 \blank{1cm}。
\choicesFour{平行}{垂直相交}{$l$ 在 $ \pi $ 上}{相交但不垂直}

\begin{redenv}
\analysisrule
\subsection*{逐步解析}

\begin{detailedanalysis}
\item \textbf{详细求解过程}
\begin{enumerate}[label=(\arabic*)]
\item \textbf{向量提取}：
直线方向向量 $\mathbf{s} = \{-2, -7, 3\}$。
平面法向量 $\mathbf{n} = \{4, -2, -2\}$。
\item \textbf{计算数量积}：
$\mathbf{s} \cdot \mathbf{n} = (-2)(4) + (-7)(-2) + 3(-2) = -8 + 14 - 6 = 0$。
积为0，说明向量垂直，即直线与平面\textbf{平行}或在平面\textbf{上}。
\item \textbf{验证点}：
取直线上一点 $P(-3, -4, 0)$。
代入平面方程：$4(-3) - 2(-4) - 0 - 3 = -12 + 8 - 3 = -7 \neq 0$。
点不在平面上，故排除“在平面上”。
\end{enumerate}

\item \textbf{最终答案}
\textbf{答案}：A
\end{detailedanalysis}
\end{redenv}

% ========== 第4题 ==========
\problemrule
\section*{【题目4】向量运算}
已知向量 $\mathbf{a}$ 垂直于向量 $ \mathbf{b} = 2\mathbf{i} - 3\mathbf{j} + \mathbf{k} $ 和 $ \mathbf{c} = \mathbf{i} - 2\mathbf{j} + 3\mathbf{k} $ ，且满足于 $ \mathbf{a} \cdot (\mathbf{i} + 2\mathbf{j} - 7\mathbf{k}) = 10 $ ，求 $\mathbf{a} =$ \blank{1cm}。
\choicesTwo
{$ -7\mathbf{i} - 5\mathbf{j} - \mathbf{k} $}
{$ 7\mathbf{i} + 5\mathbf{j} + \mathbf{k} $}
{$ -5\mathbf{i} - 3\mathbf{j} - \mathbf{k} $}
{$ 5\mathbf{i} + 3\mathbf{j} + \mathbf{k} $}

\begin{redenv}
\analysisrule
\subsection*{逐步解析}

\begin{detailedanalysis}
\item \textbf{详细求解过程}
\begin{enumerate}[label=(\arabic*)]
\item \textbf{确定方向}：
$\mathbf{a} \perp \mathbf{b}$ 且 $\mathbf{a} \perp \mathbf{c}$，故 $\mathbf{a} \parallel \mathbf{b} \times \mathbf{c}$。
$\mathbf{b} \times \mathbf{c} = \begin{vmatrix} \mathbf{i} & \mathbf{j} & \mathbf{k} \\ 2 & -3 & 1 \\ 1 & -2 & 3 \end{vmatrix} = (-9+2)\mathbf{i} - (6-1)\mathbf{j} + (-4+3)\mathbf{k} = \{-7, -5, -1\}$。
设 $\mathbf{a} = k\{-7, -5, -1\}$。
\item \textbf{利用数量积}：
$\mathbf{a} \cdot \{1, 2, -7\} = k(-7-10+7) = -10k$。
已知积为 10，所以 $-10k = 10 \implies k = -1$。
\item \textbf{求解 $\mathbf{a}$}：
$\mathbf{a} = -1 \{-7, -5, -1\} = \{7, 5, 1\}$。
\end{enumerate}

\item \textbf{最终答案}
\textbf{答案}：B
\end{detailedanalysis}
\end{redenv}

% ========== 第5题 ==========
\problemrule
\section*{【题目5】方向余弦}
向量 $\mathbf{a}$ 与 $x$ 轴与 $y$ 轴构成等角，与 $z$ 轴夹角是前者的 2 倍，下面哪一个代表的是 $\mathbf{a}$ 的方向 \blank{1cm}。
\choicesTwo
{$ \alpha=\frac{\pi}{2},\beta=\frac{\pi}{2},\gamma=\frac{\pi}{4} $}
{$ \alpha=\frac{\pi}{4},\beta=\frac{\pi}{4},\gamma=\frac{\pi}{8} $}
{$ \alpha=\frac{\pi}{4},\beta=\frac{\pi}{4},\gamma=\frac{\pi}{2} $}
{$ \alpha=\pi,\beta=\frac{\pi}{2},\gamma=\frac{\pi}{2} $}

\begin{redenv}
\analysisrule
\subsection*{逐步解析}

\begin{detailedanalysis}
\item \textbf{详细求解过程}
设 $\alpha = \beta$，$\gamma = 2\alpha$。
方向余弦性质：$\cos^2 \alpha + \cos^2 \beta + \cos^2 \gamma = 1$。
$2\cos^2 \alpha + \cos^2 2\alpha = 1$。
$2\cos^2 \alpha + (2\cos^2 \alpha - 1)^2 = 1$。
令 $u = \cos^2 \alpha$。
$2u + (2u-1)^2 = 1 \implies 2u + 4u^2 - 4u + 1 = 1 \implies 4u^2 - 2u = 0$。
$2u(2u-1) = 0$。
$u=0$ 或 $u=1/2$。
\begin{itemize}
    \item 若 $u=0$，$\alpha = \pi/2$，则 $\gamma = \pi$。选项 D 不符（D中 $\gamma=\pi/2$）。
    \item 若 $u=1/2$，$\cos \alpha = \pm \frac{\sqrt{2}}{2}$，$\alpha = \pi/4$ 或 $3\pi/4$。
    此时 $\gamma = 2\alpha = \pi/2$ 或 $3\pi/2$。
\end{itemize}
选项 C 符合 $\alpha=\pi/4, \beta=\pi/4, \gamma=\pi/2$。

\item \textbf{最终答案}
\textbf{答案}：C
\end{detailedanalysis}
\end{redenv}

% ========== 第6题 ==========
\problemrule
\section*{【题目6】单位向量}
设已知两点 $A(4,0,5)$ 和 $B(7,1,3)$，方向和 $ \overrightarrow{AB} $ 一致的单位向量为 \blank{2cm}。

\begin{redenv}
\analysisrule
\subsection*{逐步解析}

\begin{detailedanalysis}
\item \textbf{详细求解过程}
\begin{enumerate}[label=(\arabic*)]
\item \textbf{求向量 $\overrightarrow{AB}$}：
$\overrightarrow{AB} = \{7-4, 1-0, 3-5\} = \{3, 1, -2\}$。
\item \textbf{求模长}：
$|\overrightarrow{AB}| = \sqrt{3^2 + 1^2 + (-2)^2} = \sqrt{9+1+4} = \sqrt{14}$。
\item \textbf{单位化}：
$\mathbf{e} = \frac{1}{\sqrt{14}} \{3, 1, -2\} = \{\frac{3}{\sqrt{14}}, \frac{1}{\sqrt{14}}, -\frac{2}{\sqrt{14}}\}$。
\end{enumerate}

\item \textbf{最终答案}
\textbf{答案}：$ \{\frac{3}{\sqrt{14}}, \frac{1}{\sqrt{14}}, -\frac{2}{\sqrt{14}}\} $
\end{detailedanalysis}
\end{redenv}

% ========== 第7题 ==========
\problemrule
\section*{【题目7】垂直向量求参}
已知 $ \mathbf{a} = \{2, 1, 2\} $ , $ \mathbf{b} = \{4, -1, 10\} $ , $ \mathbf{c} = \mathbf{b} - \lambda \mathbf{a} $ , 且 $ \mathbf{a} \perp \mathbf{c} $ , 则 $ \lambda = $ \blank{2cm}。

\begin{redenv}
\analysisrule
\subsection*{逐步解析}

\begin{detailedanalysis}
\item \textbf{详细求解过程}
\begin{enumerate}[label=(\arabic*)]
\item \textbf{垂直条件}：
$\mathbf{a} \cdot \mathbf{c} = 0 \implies \mathbf{a} \cdot (\mathbf{b} - \lambda \mathbf{a}) = 0$。
$\mathbf{a} \cdot \mathbf{b} - \lambda |\mathbf{a}|^2 = 0 \implies \lambda = \frac{\mathbf{a} \cdot \mathbf{b}}{|\mathbf{a}|^2}$。
\item \textbf{计算}：
$\mathbf{a} \cdot \mathbf{b} = 2(4) + 1(-1) + 2(10) = 8 - 1 + 20 = 27$。
$|\mathbf{a}|^2 = 2^2 + 1^2 + 2^2 = 4 + 1 + 4 = 9$。
\item \textbf{求解}：
$\lambda = \frac{27}{9} = 3$。
\end{enumerate}

\item \textbf{最终答案}
\textbf{答案}：3
\end{detailedanalysis}
\end{redenv}

% ========== 第8题 ==========
\problemrule
\section*{【题目8】平面方程 (点法式)}
过点 $ (3,2,1) $ 与直线 $ x = y = z $ 垂直的平面方程为 \blank{2cm}。

\begin{redenv}
\analysisrule
\subsection*{逐步解析}

\begin{detailedanalysis}
\item \textbf{详细求解过程}
\begin{enumerate}[label=(\arabic*)]
\item \textbf{法向量}：
平面垂直于直线，故平面的法向量即为直线的方向向量。
直线 $x=y=z$ 的方向向量 $\mathbf{s} = \{1, 1, 1\}$。
所以 $\mathbf{n} = \{1, 1, 1\}$。
\item \textbf{点法式方程}：
$1(x-3) + 1(y-2) + 1(z-1) = 0$。
$x + y + z - 6 = 0$。
\end{enumerate}

\item \textbf{最终答案}
\textbf{答案}：$ x+y+z-6=0 $
\end{detailedanalysis}
\end{redenv}

% ========== 第9题 ==========
\problemrule
\section*{【题目9】平面方程}
过点 $ (3,-2,2) $ 与直线 $ x=\frac{y}{2}=z $ 垂直的平面方程是 \blank{2cm}。

\begin{redenv}
\analysisrule
\subsection*{逐步解析}

\begin{detailedanalysis}
\item \textbf{详细求解过程}
\begin{enumerate}[label=(\arabic*)]
\item \textbf{法向量}：
直线方向向量 $\mathbf{s} = \{1, 2, 1\}$，即为平面法向量 $\mathbf{n}$。
\item \textbf{方程}：
$1(x-3) + 2(y+2) + 1(z-2) = 0$。
$x + 2y + z - (3 - 4 + 2) = 0$。
$x + 2y + z - 1 = 0$。
\end{enumerate}

\item \textbf{最终答案}
\textbf{答案}：$ x+2y+z-1=0 $
\end{detailedanalysis}
\end{redenv}

% ========== 第10题 ==========
\problemrule
\section*{【题目10】直线方程 (点向式)}
过点 $ (2,2,3) $ 且与平面 $ 3x + y = z + 2 $ 垂直的直线方程是 \blank{2cm}。

\begin{redenv}
\analysisrule
\subsection*{逐步解析}

\begin{detailedanalysis}
\item \textbf{详细求解过程}
\begin{enumerate}[label=(\arabic*)]
\item \textbf{方向向量}：
直线垂直于平面，故直线方向向量 $\mathbf{s}$ 平行于平面法向量 $\mathbf{n}$。
平面方程 $3x + y - z - 2 = 0$，法向量 $\mathbf{n} = \{3, 1, -1\}$。
取 $\mathbf{s} = \{3, 1, -1\}$。
\item \textbf{点向式}：
$\frac{x-2}{3} = \frac{y-2}{1} = \frac{z-3}{-1}$。
\end{enumerate}

\item \textbf{最终答案}
\textbf{答案}：$ \frac{x-2}{3} = \frac{y-2}{1} = \frac{z-3}{-1} $
\end{detailedanalysis}
\end{redenv}

% ========== 第11题 ==========
\problemrule
\section*{【题目11】线面位置判定}
直线 $ \frac{x-1}{2} = \frac{y+2}{5} = \frac{z-4}{-3} $ 的方向向量 $s=$ \blank{1.5cm}，与平面 $ 2x + 5y - 3z - 4 = 0 $ 是 \blank{1.5cm} 的。

\begin{redenv}
\analysisrule
\subsection*{逐步解析}

\begin{detailedanalysis}
\item \textbf{详细求解过程}
\begin{enumerate}[label=(\arabic*)]
\item \textbf{方向向量}：$\mathbf{s} = \{2, 5, -3\}$。
\item \textbf{平面法向量}：$\mathbf{n} = \{2, 5, -3\}$。
\item \textbf{关系判别}：
$\mathbf{s}$ 与 $\mathbf{n}$ 为同一向量（或成比例），说明直线与平面法线平行，即直线与平面\textbf{垂直}。
\end{enumerate}

\item \textbf{最终答案}
\textbf{答案}：$ \{2, 5, -3\} $；垂直
\end{detailedanalysis}
\end{redenv}

% ========== 第12题 ==========
\problemrule
\section*{【题目12】平面方程}
求平行于 $y$ 轴且过点 $ P(1,5,1) $ , $ Q(3,2,-1) $ 的平面方程。

\begin{redenv}
\analysisrule
\subsection*{逐步解析}

\begin{detailedanalysis}
\item \textbf{详细求解过程}
\begin{enumerate}[label=(\arabic*)]
\item \textbf{法向量求解}：
平行于 $y$ 轴 $\implies \mathbf{n} \perp \mathbf{j}$ ($y$轴方向)。
过 $P, Q \implies \mathbf{n} \perp \overrightarrow{PQ}$。
$\overrightarrow{PQ} = \{2, -3, -2\}$。
$\mathbf{n} = \mathbf{j} \times \overrightarrow{PQ} = \begin{vmatrix} \mathbf{i} & \mathbf{j} & \mathbf{k} \\ 0 & 1 & 0 \\ 2 & -3 & -2 \end{vmatrix} = -2\mathbf{i} + 0\mathbf{j} - 2\mathbf{k} = \{-2, 0, -2\}$。
可取 $\mathbf{n} = \{1, 0, 1\}$。
\item \textbf{方程}：
$1(x-1) + 0(y-5) + 1(z-1) = 0 \implies x + z - 2 = 0$。
\end{enumerate}

\item \textbf{最终答案}
\textbf{答案}：$ x+z-2=0 $
\end{detailedanalysis}
\end{redenv}

% ========== 第13题 ==========
\problemrule
\section*{【题目13】面与线综合}
已知 $ \triangle ABC $ 的三个顶点的坐标分别为 $ A(1,2,3) $ , $ B(3,4,5) $ , $ C(2,4,7) $ ,
求：(1) $ \triangle ABC $ 所在的平面方程；
(2) $AB$ 边上的中线 $CD$ 所在的直线方程。

\begin{redenv}
\analysisrule
\subsection*{逐步解析}

\begin{detailedanalysis}
\item \textbf{详细求解过程}
\begin{enumerate}[label=(\arabic*)]
\item \textbf{求平面方程}：
$\overrightarrow{AB} = \{2, 2, 2\}$，$\overrightarrow{AC} = \{1, 2, 4\}$。
$\mathbf{n} = \overrightarrow{AB} \times \overrightarrow{AC} = \begin{vmatrix} \mathbf{i} & \mathbf{j} & \mathbf{k} \\ 2 & 2 & 2 \\ 1 & 2 & 4 \end{vmatrix} = 4\mathbf{i} - 6\mathbf{j} + 2\mathbf{k}$。
取 $\mathbf{n} = \{2, -3, 1\}$。
方程：$2(x-1) - 3(y-2) + 1(z-3) = 0 \implies 2x - 3y + z + 1 = 0$。
\item \textbf{求直线方程}：
$D$ 为 $AB$ 中点：$D(\frac{1+3}{2}, \frac{2+4}{2}, \frac{3+5}{2}) = (2, 3, 4)$。
直线 $CD$ 过 $C(2, 4, 7)$ 和 $D(2, 3, 4)$。
方向向量 $\overrightarrow{CD} = \{0, -1, -3\}$。
方程：$\frac{x-2}{0} = \frac{y-3}{-1} = \frac{z-4}{-3}$（参数式为 $x=2, \frac{y-3}{1} = \frac{z-4}{3}$）。
\end{enumerate}

\item \textbf{最终答案}
\textbf{答案}：(1) $ 2x-3y+z+1=0 $; (2) $ \begin{cases} x=2 \\ 3y-z-5=0 \end{cases} $
\end{detailedanalysis}
\end{redenv}

% ========== 第14题 ==========
\problemrule
\section*{【题目14】直线方程}
求过点 $ (-1,-4,3) $ 并与两直线 $ L_{1}\begin{cases}2x-4y+z=1\\ x+3y+5=0\end{cases} $ 和 $ L_{2}\begin{cases}x=2+4t\\ y=-1-t\\ z=-3+2t\end{cases} $ 都垂直的直线方程。

\begin{redenv}
\analysisrule
\subsection*{逐步解析}

\begin{detailedanalysis}
\item \textbf{详细求解过程}
\begin{enumerate}[label=(\arabic*)]
\item \textbf{求 $L_1$ 方向}：
$\mathbf{s}_1 = \mathbf{n}_1 \times \mathbf{n}_2 = \{2, -4, 1\} \times \{1, 3, 0\} = \{-3, 1, 10\}$。
\item \textbf{求 $L_2$ 方向}：
$\mathbf{s}_2 = \{4, -1, 2\}$。
\item \textbf{求所求直线方向 $\mathbf{s}$}：
$\mathbf{s} = \mathbf{s}_1 \times \mathbf{s}_2 = \begin{vmatrix} \mathbf{i} & \mathbf{j} & \mathbf{k} \\ -3 & 1 & 10 \\ 4 & -1 & 2 \end{vmatrix} = 12\mathbf{i} - (-46)\mathbf{j} + (-1)\mathbf{k} = \{12, 46, -1\}$。
\item \textbf{写方程}：
$\frac{x+1}{12} = \frac{y+4}{46} = \frac{z-3}{-1}$。
\end{enumerate}

\item \textbf{最终答案}
\textbf{答案}：$ \frac{x+1}{12} = \frac{y+4}{46} = \frac{z-3}{-1} $
\end{detailedanalysis}
\end{redenv}

% ========== 第15题 ==========
\problemrule
\section*{【题目15】平面方程}
求通过点 $ M_{1}(3,-5,1) $ ， $ M_{2}(4,1,2) $ 且垂直于平面 $ x-8y+3z-1=0 $ 的平面方程。

\begin{redenv}
\analysisrule
\subsection*{逐步解析}

\begin{detailedanalysis}
\item \textbf{详细求解过程}
\begin{enumerate}[label=(\arabic*)]
\item \textbf{法向量求解}：
法向量 $\mathbf{n}$ 垂直于平面法向 $\mathbf{n}_0 = \{1, -8, 3\}$，且垂直于 $\overrightarrow{M_1 M_2} = \{1, 6, 1\}$。
$\mathbf{n} = \mathbf{n}_0 \times \overrightarrow{M_1 M_2} = \begin{vmatrix} \mathbf{i} & \mathbf{j} & \mathbf{k} \\ 1 & -8 & 3 \\ 1 & 6 & 1 \end{vmatrix} = (-8-18)\mathbf{i} - (1-3)\mathbf{j} + (6+8)\mathbf{k} = \{-26, 2, 14\}$。
取 $\mathbf{n} = \{-13, 1, 7\}$。
\item \textbf{方程}：
$-13(x-3) + 1(y+5) + 7(z-1) = 0$。
$-13x + y + 7z - (-39 - 5 + 7) = 0$。
$-13x + y + 7z + 37 = 0$（或 $13x - y - 7z - 37 = 0$）。
\end{enumerate}

\item \textbf{最终答案}
\textbf{答案}：$ 13x - y - 7z - 37 = 0 $
\end{detailedanalysis}
\end{redenv}

\end{document}
